% =====================================================================
\section{The closure-response operator}\label{sec:definition}
% =====================================================================

\subsection{Definition}

Let $M$ be a closure substrate: a connected finite undirected graph
$M=(V,E)$, a finite simplicial complex with chosen Laplacian, or a
projected continuum coordinate. Let $L_M$ be the corresponding
Laplacian (graph Laplacian $L = D - A$, simplicial Laplacian, or
continuum operator $-\Delta$ with chosen boundary conditions).
Let $\Ph = (1+\sqrt{5})/2$ be the golden ratio, with $\Ph^{-1} = \Ph - 1$
and $\Ph^{-2} = 2 - \Ph \approx 0.381966$.

The \emph{closure-response operator} is
\begin{equation}\label{eq:cphi}
\Cph \;=\; L_M + \Ph^{-2} I.
\end{equation}
For a non-negative localised source $f$ on $M$, the
\emph{closure response field} is
\begin{equation}\label{eq:psi}
\psi \;=\; \Cph^{-1} f \;=\; (L_M + \Ph^{-2} I)^{-1} f.
\end{equation}

\subsection{Hypotheses on $(M, L_M)$}\label{ssec:hypotheses}

The properties developed in \S\ref{ssec:positivity}--\S\ref{ssec:opnorm}
require:

\begin{itemize}\itemsep=2pt
\item \textbf{(H1) Self-adjointness.} $L_M$ is self-adjoint on the
  $L^{2}$ inner product on $M$ (with counting measure on a finite
  graph, with Lebesgue measure with appropriate boundary conditions
  in the continuum case).
\item \textbf{(H2) Non-negativity.} $L_M \geq 0$ as a
  quadratic form: $\langle f, L_M f\rangle \geq 0$ for all $f$.
\item \textbf{(H3) Known spectral bottom.} On a finite graph, $M$
  is connected (so the kernel of $L_M$ is exactly the constant
  vector and $\lambda_{\min}(L_M) = 0$). In the continuum case,
  the spectral bottom $\inf \sigma(L_M)$ is explicitly known on
  the chosen domain (it equals $0$ for the full line; equals the
  first Dirichlet eigenvalue $> 0$ on a bounded Dirichlet
  interval). The full-line case has $\inf \sigma(L_M) = 0$ as
  spectral bottom but not as an eigenvalue, which is sufficient
  for the operator-norm identity in
  \S\ref{ssec:opnorm}; H3 does not require a finite-dimensional
  zero eigenspace.
\end{itemize}

Three concrete settings illustrate the hypothesis class:
\begin{itemize}\itemsep=2pt
\item the standard combinatorial Laplacian on a connected finite
  undirected graph (the 600-cell case, $\lambda_{\min}(L_M) = 0$);
\item the continuum $L = -d^{2}/dx^{2}$ on the full line with
  decay-at-infinity (spectral bottom $\inf \sigma(L_M) = 0$,
  not attained as an eigenvalue; used for the closed-form Green's
  function in \S\ref{ssec:continuum});
\item the continuum $L = -d^{2}/dx^{2}$ on a bounded interval
  with Dirichlet boundary conditions ($\lambda_{\min}(L_M) > 0$).
\end{itemize}
Substrates outside this class — projected coordinates with
non-standard boundaries, weighted Laplacians whose weight function
is unbounded, or operators with negative spectrum — require their
own analysis, which we do not give here.

\subsection{Positive definiteness}\label{ssec:positivity}

Under (H1)--(H3) on a finite connected graph, $L_M$ has a smallest
eigenvalue $\lambda_{\min}(L_M) = 0$ with one-dimensional
eigenspace (the constant vector). For $\Cph = L_M + \Ph^{-2} I$,
\[
\lambda_{\min}(\Cph) \;=\; \lambda_{\min}(L_M) + \Ph^{-2}
                    \;=\; \Ph^{-2} \;>\; 0,
\]
so $\Cph$ is strictly positive definite and $\Cph^{-1}$ is
well-defined and bounded.

\subsection{Operator-norm bound}\label{ssec:opnorm}

The operator norm of $\Cph^{-1}$ is the reciprocal of the
spectral bottom of $\Cph$:
\begin{equation}\label{eq:opnorm}
\|\Cph^{-1}\| \;=\; 1/\inf \sigma(\Cph)
              \;=\; 1/(\inf \sigma(L_M) + \Ph^{-2}).
\end{equation}
On finite graphs the spectral bottom is the smallest eigenvalue
$\lambda_{\min}(L_M)$ (and is attained); on the full-line
continuum case (\S\ref{ssec:continuum}) it is $\inf \sigma(L_M)=0$
as spectral bottom but \emph{not} as an eigenvalue. On any
substrate where $\inf \sigma(L_M) = 0$ (e.g.\ the connected finite
graph $\Rsixhundred$, or the full-line continuum case), this
reduces to the identity
\begin{equation}\label{eq:opnorm_zero_mode}
\|\Cph^{-1}\| \;=\; \Ph^{2} \;\approx\; 2.618034.
\end{equation}
On substrates where $\inf \sigma(L_M) > 0$ (e.g.\ Dirichlet
intervals), Eq.~\eqref{eq:opnorm} gives the strict inequality
$\|\Cph^{-1}\| < \Ph^{2}$. The response-amplification ceiling is,
in either case, $\|\psi\|_{2} \leq \|\Cph^{-1}\|\, \|f\|_{2}$. On
substrates with $\inf \sigma(L_M) > 0$ a finite ceiling is already
guaranteed by the Dirichlet structure alone; on substrates with
$\inf \sigma(L_M) = 0$ (the cases of interest here) the shift
$\Ph^{-2}$ is what guarantees a finite ceiling and pins it at
$\Ph^{2}$. Numerically reproduced as $\|\Cph^{-1}\| =
2.618034$ on $\Rsixhundred$ (\texttt{repro/verify\_kernel.py},
\S\ref{ssec:opnorm_check}); this matches the closed-form $\Ph^{2}$
to machine precision.

\subsection{Continuum projection}\label{ssec:continuum}

In one projected coordinate $x \in \mathbb{R}$ with
$L_{\Ph} = -d^{2}/dx^{2} + \Ph^{-2}$, the Green's function
$G(x)$ satisfies $L_{\Ph} G = \delta_{0}$ and is the closed-form
exponential
\begin{equation}\label{eq:green_continuum}
G(x) \;=\; \frac{\Ph}{2}\, e^{-|x|/\Ph}.
\end{equation}
The decay scale is $\Ph$ — the same constant that appears in the
shift, by construction. Normalised, the kernel is
$\kappa(x) = e^{-|x|/\Ph}$ with unit value at the source.

This continuum Green's function is the comparison object for the
discrete-to-continuum agreement test (\S\ref{sec:agreement}):
the discrete response $\psi(v) = \Cph^{-1} f(v)$ at a vertex $v$ at
chord distance $\|v - v_{\mathrm{src}}\|$ from a localised source
is compared to $G(\|v - v_{\mathrm{src}}\|)$.

\subsection{Disclosure: $\Ph^{-2}$ is a design-level shift}

The shift $\Ph^{-2}$ is chosen so that:
\begin{enumerate}\itemsep=2pt
\item $\Cph$ is strictly positive definite (the smallest eigenvalue
  is exactly $\Ph^{-2}$);
\item $\Ph^{2}$ is the zero-bottom operator-norm
  \emph{identity} ($\|\Cph^{-1}\| = \Ph^{2}$ when
  $\inf\sigma(L_M)=0$) and the general response-amplification
  \emph{ceiling} ($\|\Cph^{-1}\| \leq \Ph^{2}$ on positive-bottom
  substrates), and the continuum decay scale is $\Ph$
  (Eq.~\eqref{eq:green_continuum}). Zero-bottom identity, general
  ceiling, and continuum decay scale are all fixed by the single
  design choice $\Ph^{-2}$, giving one dimensional parameter
  throughout the operator;
\item the continuum projection (Eq.~\eqref{eq:green_continuum})
  has decay scale $\Ph$, not a free length parameter.
\end{enumerate}
We do \emph{not} derive $\Ph^{-2}$ from a closure functional or
symmetry argument. It is a design-level choice motivated by
(1)--(3); we report this explicitly and treat formal derivation as
an open build (\S\ref{sec:limitations}). The companion
adaptive-closure-transport
preprint~\citep{SmartAdaptiveClosureTransport2026} formulates the
selection-layer dynamics over $W$-space that would, if delivered,
constrain the shift further; that derivation is not delivered
there or here.
